\section{Attention \& Transformers}
Early sequence processing models relied primarily on recurrent neural network (RNN) encoder-decoder architectures. 
In these models, the encoder compresses the input sequence into a fixed-size representation, which the decoder then uses to generate the output. 
However, this fixed bottleneck limits performance. The introduction of attention mechanisms demonstrated that creating a dynamic, context-aware representation of the input sequence significantly improves results. 

Building on the success of attention, researchers sought to eliminate the sequential bottleneck of RNNs entirely. 
This led to the development of more flexible architectures capable of processing sequences in parallel while still effectively capturing long-range dependencies.

\subsection{Convolutional Sequence Models}
Convolutional Neural Networks (CNNs) offer a natural alternative for sequence processing. 
By using convolutional filters, they efficiently capture local dependencies, and stacking these layers allows the model to expand its receptive field to capture longer-range relationships.

CNNs provide two primary advantages over RNNs:
\begin{itemize}
    \item \textbf{Parallelization:} Unlike RNNs, which must process tokens sequentially, CNNs process the entire sequence hierarchically. 
    This enables massive parallelization and significantly faster training times.
    
    \item \textbf{Parameter Efficiency:} The parameter-sharing nature of convolutional filters makes CNNs highly efficient, 
    reducing both the total number of parameters and overall computational cost.
\end{itemize}

Despite these strengths, traditional CNNs are constrained by a fixed input size. This limitation complicates their application 
to variable-length sequences, often requiring workarounds that make them less flexible in certain scenarios.

\paragraph{Dilated Convolutions}
To capture long-range dependencies without inflating the parameter count, CNNs frequently employ \textbf{dilated convolutions}. 
Unlike standard convolutions, a dilated filter is applied over a broader area by introducing gaps between the sampled input values. 
The size of these intervals is determined by the \textit{dilation rate}. This technique effectively expands the network's receptive field while preserving both parameter efficiency and computational tractability.

For instance, consider a standard $3 \times 3$ convolutional filter applied to adjacent pixels. This configuration corresponds to a dilation rate of $1$, yielding a $3 \times 3$ receptive field. 
If the dilation rate is increased to $2$, the filter skips one pixel between each consecutive sample. Consequently, the filter's effective receptive field expands to a $5 \times 5$ area, despite still relying on the same $9$ trainable parameters.